\section{Rangkuman}
\begin{itemize}
    \item DES adalah nama standar, sedangkan algoritmanya disebut DEA. DES
    termasuk cipher blok kunci-simetri dengan blok plainteks 64 bit. Kunci
    eksternalnya 64 bit, tetapi hanya 56 bit yang dipakai; delapan bit paritas
    pada posisi 8, 16, 24, 32, 40, 48, 56, dan 64 tidak dipakai.
    \item Pada tahun 1972 NBS (sekarang NIST) meminta standar cipher blok.
    Tahun 1974 Horst Feistel di IBM merancang Lucifer dengan kunci 128 bit dan
    blok 128 bit. Lucifer berkembang menjadi DES dengan kunci 56 bit, blok
    64 bit, dan 16 putaran, lalu disetujui NBS pada tahun 1976 setelah dinilai
    NSA.
    \item Alur enkripsi: permutasi awal (IP), 16 putaran \textit{enciphering},
    lalu invers permutasi awal ($\text{IP}^{-1}$) atas gabungan blok kiri dan
    blok kanan.
    \item Fungsi putaran $f$ memperluas blok kanan 32 bit menjadi 48 bit dengan
    tabel $E$, di-XOR dengan kunci internal 48 bit menjadi blok $A$, dipecah
    menjadi delapan kelompok 6 bit, disubstitusi melalui delapan kotak-S
    (6 bit menjadi 4 bit) menjadi blok $B$ 32 bit, lalu dipermutasi oleh $P$.
    \item Pada setiap kotak-S, bit pertama dan keenam masukan menentukan nomor
    baris (0--3), sedangkan bit kedua sampai kelima menentukan nomor kolom
    (0--15). Setiap baris adalah permutasi $0, 1, \dots, 15$. Mengubah satu bit
    masukan mengubah paling sedikit dua bit luaran, yang mendukung prinsip
    konfusi.
    \item Pembangkitan kunci: PC-1 membuang bit paritas dan membagi kunci
    menjadi $C_0$ dan $D_0$ (masing-masing 28 bit), lalu pergeseran siklik
    menurut jadwal $1,1,2,2,\dots,1$ dan PC-2 memilih 48 bit menjadi $K_i$.
    \item Dekripsi memakai algoritma yang sama dengan urutan kunci terbalik
    $K_{16}, \dots, K_1$.
    \item Ruang kunci DES berisi $2^{56} = 72.057.594.037.927.936$ kunci. Pada
    laju satu juta kunci per detik, seluruh ruang kunci memerlukan waktu sekitar
    1142 tahun, sedangkan kunci diharapkan ditemukan sebelum $2^{55}$ iterasi.
    \item DES dipecahkan secara terbuka: kontes RSA Security 1997 dengan hadiah
    \$10.000 dimenangkan DESCHALL, mesin EFF senilai \$250.000 berisi 1.856 chip
    pada 1998 menemukan kunci dalam sedikit lebih dari dua hari, dan gabungan
    mesin itu dengan lebih dari 100.000 komputer pada 1999 menemukan kunci
    kurang dari satu hari.
    \item Terdapat empat kunci lemah yang membuat seluruh 16 kunci internal
    sama, sehingga $E_K(E_K(x)) = x$, yaitu
    \begin{center}
    \small
    \code{0101010101010101} \quad \code{1F1F1F1F0E0E0E0E} \quad
    \code{E0E0E0E0F1F1F1F1} \quad \code{FEFEFEFEFEFEFEFE}.
    \end{center}
    \item Efek longsoran DES kuat: pembalikan satu bit plainteks pada vektor uji
    resmi mengubah 33 dari 64 bit cipherteks.
    \item DES dapat dioperasikan dengan mode ECB, CBC, OFB, CFB, dan mode
    \textit{counter}; mode ECB paling sering dipakai pada paket komersial dan
    dikenal sebagai DES-ECB.
    \item Implementasi perangkat keras DES dapat mengenkripsikan 16,8 juta blok
    per detik (sekitar satu gigabit per detik), sedangkan implementasi perangkat
    lunak pada IBM 3090 mencapai 32.000 blok per detik.
    \item Double DES memakai dua kunci, $C = E_{K_2}(E_{K_1}(P))$, dan tumbang
    oleh serangan \textit{meet-in-the-middle}: $X = E_{K_1}(P) = D_{K_2}(C)$.
    Penyerang menyimpan $2^{56}$ nilai $X$ dari semua $K_1$, lalu membandingkannya
    dengan hasil dekripsi $C$ memakai semua $K_2$, dan menguji kecocokan dengan
    pasangan plainteks-cipherteks lain.
    \item Triple DES memakai DES tiga kali untuk mencegah serangan
    \textit{meet-in-the-middle}. Mode EEE dapat disederhanakan menjadi mode EDE
    agar kompatibel dengan DES tunggal ketika $K_1=K_2=K_3$. Versi dua kunci
    berukuran $56+56=112$ bit, sedangkan versi tiga kunci berukuran
    $56+56+56=168$ bit.
\end{itemize}
